%=====================================================================
% §10 BENCHMARKS & EVALUATION  (cross-cutting)
% MASTER §8  (#94-105)      OWNER: Systems lead (Profile C) / entry (Profile D) | TARGET: ~2 pp
% CLOSE WITH:
%   (Q1) Are current benchmarks adequate to measure MARL-specific robustness?
%   (Q2) Open gap? (standardized robustness metrics; agent-failure & communication-attack benches)
%=====================================================================
\section{Benchmarks and Evaluation}
\label{sec:benchmarks}

\subsection{Robustness-Oriented Benchmarks and Testbeds}
% \cite{} Robust Gymnasium %TODO; \cite{li2025starcraft+} (adversary paradigm);
% \cite{barta2025measuring} (partial agent failure, x-ref §7); general MARL: SMACv2 %TODO, MATE %TODO, Mava %TODO.
TODO.

\subsection{Robustness Testing and Empirical Studies}
% \cite{} robustness testing on critical agents %TODO (x-ref §2,§5); comprehensive testing %TODO;
% empirical study on robustness & resilience %TODO 2510.11824 (x-ref §7);
% observation poisoning %TODO; domain studies (traffic signal, urban driving, x-ref §11).
TODO.

\begin{table}[t]\centering\footnotesize
\caption{Benchmarks/testbeds for robust MARL (MASTER \S8).}
\label{tab:bench}
\begin{tabular}{lllll}
\toprule
Benchmark & Perturbation types & Metric & Agents/scale & Open-source \\
\midrule
TODO & & & & \\
\bottomrule
\end{tabular}
\end{table}

% [Closing] answer Q1 & Q2 -- motivates a "standardized robustness benchmark" challenge in §12.
